
{\actuality} 
% Обзор, введение в тему, обозначение места данной работы в
% мировых исследованиях и~т.\:п., можно использовать ссылки на~другие
% работы~\autocite{Gosele1999161,Lermontov}
% (если их~нет, то~в~автореферате
% автоматически пропадёт раздел <<Список литературы>>). Внимание! Ссылки
% на~другие работы в~разделе общей характеристики работы можно
% использовать только при использовании \verb!biblatex! (из-за технических
% ограничений \verb!bibtex8!. Это связано с тем, что одна
% и~та~же~характеристика используются и~в~тексте диссертации, и в
% автореферате. В~последнем, согласно ГОСТ, должен присутствовать список
% работ автора по~теме диссертации, а~\verb!bibtex8! не~умеет выводить в~одном
% файле два списка литературы).
% При использовании \verb!biblatex! возможно использование исключительно
% в~автореферате подстрочных ссылок
% для других работ командой \verb!\autocite!, а~также цитирование
% собственных работ командой \verb!\cite!. Для этого в~файле
% \verb!common/setup.tex! необходимо присвоить положительное значение
% счётчику \verb!\setcounter{usefootcite}{1}!.

% \ifsynopsis
% Этот абзац появляется только в~автореферате.
% Для формирования блоков, которые будут обрабатываться только в~автореферате,
% заведена проверка условия \verb!\!\verb!ifsynopsis!.
% Значение условия задаётся в~основном файле документа (\verb!synopsis.tex! для
% автореферата).
% \else
% Этот абзац появляется только в~диссертации.
% Через проверку условия \verb!\!\verb!ifsynopsis!, задаваемого в~основном файле
% документа (\verb!dissertation.tex! для диссертации), можно сделать новую
% команду, обеспечивающую появление цитаты в~диссертации, но~не~в~автореферате.
% \fi


% ~\autocite{bulatov2022recurrent}




% {\progress}
% Этот раздел должен быть отдельным структурным элементом по
% ГОСТ, но он, как правило, включается в описание актуальности
% темы. Нужен он отдельным структурынм элемементом или нет ---
% смотрите другие диссертации вашего совета, скорее всего не нужен.

Данная диссертация посвящена ...

Постановка проблемы

Отсюда возникает проблема настоящего исследования:
как обучать языко- и доменно- специфичные модели максимально переисполь
зуя уже обученные модели.

В ряде работ...


{\aim} данной работы является исследование 


Для~достижения поставленной цели необходимо было решить следующие {\tasks}:
\begin{enumerate}[beginpenalty=10000] % https://tex.stackexchange.com/a/476052/104425
    \item Предложить подход и разработать метод дообучения 
    \item Применить разработанный метод переноса знаний для обучения языковых моделей 
    \item Применить разработанный метод переноса знаний для обучения языковых моделей разговорного домена ля русского и английского языков.
    \item Разработать и опубликовать в открытом доступе программный комплекс для дообучения языковых моделей.
    \item Исследовать эффективность полученных моделей на задачах 
\end{enumerate}


{\novelty}
\begin{enumerate}[beginpenalty=10000] % https://tex.stackexchange.com/a/476052/104425
  \item Впервые Был предложен оригинальный метод 
  \item Впервые было проведено 
  \item Было выполнено оригинальное исследование \ldots
\end{enumerate}

{\influence} 
\textbf{Теоретическая и практическая значимость.} диссертационной работы заключается в следующих положениях:
1. Предложен метод дообучения 
2. Показано, что дообучение ...позволяет получать более высокие оценки 

К практической значимости относятся следующие положения.
1. Предложенный метод переноса знаний позволяет ускорить процесс до
обучения языковых моделей.
2. С помощью предложенного метода переноса знаний обучены: языковая
модель 
3. установлены новые максимальные значения показателей качества, определяющие текущий
статус научного прогресса.
4. Обученные в рамках диссертационной работы языковые модели выло
жены в открытый доступ и могут быть использованы для улучшения
решений 
5. Обученные в рамках диссертационной работы модели 

{\methods} В ходе работы была применена методология численного эксперимента для исследования рассматриваемых в диссертации задач, применены методы теории вероятностей, машинного обучения и теории нейронных сетей. Также были применены методы разработки приложений на языке программирования Python, языке для написания скриптов Bash, программных библиотеках для машинного обучения PyTorch, Hugging Face transformers.

{\defpositions}
\begin{enumerate}[beginpenalty=10000] % https://tex.stackexchange.com/a/476052/104425
  \item Предложенная архитектура рекуррентного Трансформера с памятью (RMT) позволяет увеличить размер контекста базовой модели Трансформер. 
  \item Предложенный метод обучения RMT с расписанием повышает качество решения задач с большим контекстом. 
  \item При обучении с расписанием RMT демонстрирует способность к обобщению на новые задачи и задачи с контекстом, превышающим обучающий.
\end{enumerate}

% В папке Documents можно ознакомиться с решением совета из Томского~ГУ
% (в~файле \verb+Def_positions.pdf+), где обоснованно даются рекомендации
% по~формулировкам защищаемых положений.

{\reliability} полученных результатов обеспечивается методикой проведения численных экспериментов. Код для проведения экспериментов выложен в репозиториях на github, веса обученных моделей и наборы данных выложены в каталоги Hugging Face. Это позволяет воспроизвести результаты экспериментов, проведенных в данной работе. Результаты диссертации находятся в соответствии с результатами, полученными другими авторами для других моделей и задач.


{\probation} Основные результаты работы данной диссертации докладывались на международных конференциях и семинарах:

\begin{itemize}
    \item The Thirty-Sixth Annual Conference on Neural Information Processing Systems, November 28 - December 9 2022, New Orleans, USA

    \item Fall into ML 2023: 2nd Conference on Machine Learning, October 26-28 2023, Moscow, Russia

    \item The 2023 Conference on Empirical Methods in Natural Language Processing, December 6–10 2023, Singapore

    \item The 38th Annual AAAI Conference on Artificial Intelligence, February 20-27, 2024, Vancouver, Canada

    \item The Forty-First International Conference on Machine Learning, Next Generation of Sequence Modeling Architectures Workshop, July 26, 2024, Vienna, Austria

    \item The 62nd Annual Meeting of the Association for Computational Linguistics, TextGraphs-17 Workshop, August 15 2024, Bangkok, Thailand

    \item XXVI Международная научно-техническая конференция "Нейроинформатика-2024", 25 октября 2024, Москва, Россия    

    \item Fall into ML 2024: 3rd Conference on Machine Learning, October 25-26 2023, Moscow, Russia

    \item AI Journey Conference, 12-13 декабря 2024, Москва

    \item The Thirty-Eighth Annual Conference on Neural Information Processing Systems, Datasets and Benchmarks Track, December 10-15, 2024, Vancouver, Canada

\end{itemize}

Кроме того, модели и наборы данных, полученные в результате работы над диссертацией, опубликованы каталоге Hugging Face даи находят свое применение у ее пользователей (например, на 25 декабря 2024 г. набор данных BABILong был скачан 5832 раз в период c 26 ноября через каталог моделей Hugging Face, а репозиторий с имплементацией модели RMT имеет 759 звезд на github).

{\publications} Основные результаты по теме диссертации изложены в 7 печатных изданиях, 2 из которых изданы в периодических научных журналах, индексируемых Web of Science, 5 — в тезисах докладов, индексируемых в Scopus.

{\contribution} Автором диссертации были полностью получены результаты в работах \cite{bulatov2022recurrent}.
%!!!!!!!!!!!!!!! TODO: Add OMMN and Trudy papers
В работе \cite{bulatov2023scaling} автор провел эксперименты по моделированию естественного языка и решению синтетических задача, построил визуализации к соответствующим разделом и работал над текстом совместно с соавторами. 

Результаты, доложенные на конференции [18], полно
стью получены автором диссертации. В работе [19] автором было произведено
обучение языковых моделей предложенным методом переноса знаний и проведе
ны все эксперименты, Архипов Михаил реализовал первую версию программно
го кода для пересборки словаря, публично доступная версия пересборки словаря
реализована автором. В работе [20] автором, совместно с Петровым Максимом,
была реализована базовая модель, адаптирована базовая модель для русско
го языка, проведены эксперименты с языковыми моделями, собраны итоговые
решения для соревнования по разрешению кореференции и анафоры «Dialogue
Evaluation 2019», Ле Тхе Ань проводил эксперименты с моделью, использующей
кореферентные связи между предложений. В работе [21] автором реализова
ны модели для поиска ответа на вопрос в тексте. В работе [22] автором была
обучена предложенным методом переноса знаний модель Славянский BERT, ко
торая позволила занять первое место по двум из трех метрик на соревновании
«BSNLP 2019 Shared Task», Архиповым Михаилом и Трофимовой Марией были
реализованы модели для распознавания именованных сущностей и проведены
эксперименты с использованием модели Славянский BERT, Сорокин Алексей ре
ализовал алгоритмы для нормализации извлеченных сущностей.
Объем и структура работы. Диссертация состоит из введения, пяти
глав, заключения и двух приложений. Полный объём диссертации составляет
121 страницу, включая 31 рисунок и 19 таблиц. Список литературы содержит
135 наименований.


% \ifnumequal{\value{bibliosel}}{0}
% {%%% Встроенная реализация с загрузкой файла через движок bibtex8. (При желании, внутри можно использовать обычные ссылки, наподобие `\cite{vakbib1,vakbib2}`).
%     {\publications} Основные результаты по теме диссертации изложены
%     в~XX~печатных изданиях,
%     X из которых изданы в журналах, рекомендованных ВАК,
%     X "--- в тезисах докладов.
% }%
% {%%% Реализация пакетом biblatex через движок biber
%     \begin{refsection}[bl-author, bl-registered]
%         % Это refsection=1.
%         % Процитированные здесь работы:
%         %  * подсчитываются, для автоматического составления фразы "Основные результаты ..."
%         %  * попадают в авторскую библиографию, при usefootcite==0 и стиле `\insertbiblioauthor` или `\insertbiblioauthorgrouped`
%         %  * нумеруются там в зависимости от порядка команд `\printbibliography` в этом разделе.
%         %  * при использовании `\insertbiblioauthorgrouped`, порядок команд `\printbibliography` в нём должен быть тем же (см. biblio/biblatex.tex)
%         %
%         % Невидимый библиографический список для подсчёта количества публикаций:
%         \printbibliography[heading=nobibheading, section=1, env=countauthorvak,          keyword=biblioauthorvak]%
%         \printbibliography[heading=nobibheading, section=1, env=countauthorwos,          keyword=biblioauthorwos]%
%         \printbibliography[heading=nobibheading, section=1, env=countauthorscopus,       keyword=biblioauthorscopus]%
%         \printbibliography[heading=nobibheading, section=1, env=countauthorconf,         keyword=biblioauthorconf]%
%         \printbibliography[heading=nobibheading, section=1, env=countauthorother,        keyword=biblioauthorother]%
%         \printbibliography[heading=nobibheading, section=1, env=countregistered,         keyword=biblioregistered]%
%         \printbibliography[heading=nobibheading, section=1, env=countauthorpatent,       keyword=biblioauthorpatent]%
%         \printbibliography[heading=nobibheading, section=1, env=countauthorprogram,      keyword=biblioauthorprogram]%
%         \printbibliography[heading=nobibheading, section=1, env=countauthor,             keyword=biblioauthor]%
%         \printbibliography[heading=nobibheading, section=1, env=countauthorvakscopuswos, filter=vakscopuswos]%
%         \printbibliography[heading=nobibheading, section=1, env=countauthorscopuswos,    filter=scopuswos]%
%         %
%         \nocite{*}%
%         %
%         {\publications} Основные результаты по теме диссертации изложены в~\arabic{citeauthor}~печатных изданиях,
%         \arabic{citeauthorvak} из которых изданы в журналах, рекомендованных ВАК\sloppy%
%         \ifnum \value{citeauthorscopuswos}>0%
%             , \arabic{citeauthorscopuswos} "--- в~периодических научных журналах, индексируемых Web of~Science и Scopus\sloppy%
%         \fi%
%         \ifnum \value{citeauthorconf}>0%
%             , \arabic{citeauthorconf} "--- в~тезисах докладов.
%         \else%
%             .
%         \fi%
%         \ifnum \value{citeregistered}=1%
%             \ifnum \value{citeauthorpatent}=1%
%                 Зарегистрирован \arabic{citeauthorpatent} патент.
%             \fi%
%             \ifnum \value{citeauthorprogram}=1%
%                 Зарегистрирована \arabic{citeauthorprogram} программа для ЭВМ.
%             \fi%
%         \fi%
%         \ifnum \value{citeregistered}>1%
%             Зарегистрированы\ %
%             \ifnum \value{citeauthorpatent}>0%
%             \formbytotal{citeauthorpatent}{патент}{}{а}{}\sloppy%
%             \ifnum \value{citeauthorprogram}=0 . \else \ и~\fi%
%             \fi%
%             \ifnum \value{citeauthorprogram}>0%
%             \formbytotal{citeauthorprogram}{программ}{а}{ы}{} для ЭВМ.
%             \fi%
%         \fi%
%         % К публикациям, в которых излагаются основные научные результаты диссертации на соискание учёной
%         % степени, в рецензируемых изданиях приравниваются патенты на изобретения, патенты (свидетельства) на
%         % полезную модель, патенты на промышленный образец, патенты на селекционные достижения, свидетельства
%         % на программу для электронных вычислительных машин, базу данных, топологию интегральных микросхем,
%         % зарегистрированные в установленном порядке.(в ред. Постановления Правительства РФ от 21.04.2016 N 335)
%     \end{refsection}%
%     \begin{refsection}[bl-author, bl-registered]
%         % Это refsection=2.
%         % Процитированные здесь работы:
%         %  * попадают в авторскую библиографию, при usefootcite==0 и стиле `\insertbiblioauthorimportant`.
%         %  * ни на что не влияют в противном случае
%         \nocite{vakbib2}%vak
%         \nocite{patbib1}%patent
%         \nocite{progbib1}%program
%         \nocite{bib1}%other
%         \nocite{confbib1}%conf
%     \end{refsection}%
%         %
%         % Всё, что вне этих двух refsection, это refsection=0,
%         %  * для диссертации - это нормальные ссылки, попадающие в обычную библиографию
%         %  * для автореферата:
%         %     * при usefootcite==0, ссылка корректно сработает только для источника из `external.bib`. Для своих работ --- напечатает "[0]" (и даже Warning не вылезет).
%         %     * при usefootcite==1, ссылка сработает нормально. В авторской библиографии будут только процитированные в refsection=0 работы.
% }

% При использовании пакета \verb!biblatex! будут подсчитаны все работы, добавленные
% в файл \verb!biblio/author.bib!. Для правильного подсчёта работ в~различных
% системах цитирования требуется использовать поля:
% \begin{itemize}
%         \item \texttt{authorvak} если публикация индексирована ВАК,
%         \item \texttt{authorscopus} если публикация индексирована Scopus,
%         \item \texttt{authorwos} если публикация индексирована Web of Science,
%         \item \texttt{authorconf} для докладов конференций,
%         \item \texttt{authorpatent} для патентов,
%         \item \texttt{authorprogram} для зарегистрированных программ для ЭВМ,
%         \item \texttt{authorother} для других публикаций.
% \end{itemize}
% Для подсчёта используются счётчики:
% \begin{itemize}
%         \item \texttt{citeauthorvak} для работ, индексируемых ВАК,
%         \item \texttt{citeauthorscopus} для работ, индексируемых Scopus,
%         \item \texttt{citeauthorwos} для работ, индексируемых Web of Science,
%         \item \texttt{citeauthorvakscopuswos} для работ, индексируемых одной из трёх баз,
%         \item \texttt{citeauthorscopuswos} для работ, индексируемых Scopus или Web of~Science,
%         \item \texttt{citeauthorconf} для докладов на конференциях,
%         \item \texttt{citeauthorother} для остальных работ,
%         \item \texttt{citeauthorpatent} для патентов,
%         \item \texttt{citeauthorprogram} для зарегистрированных программ для ЭВМ,
%         \item \texttt{citeauthor} для суммарного количества работ.
% \end{itemize}
% % Счётчик \texttt{citeexternal} используется для подсчёта процитированных публикаций;
% % \texttt{citeregistered} "--- для подсчёта суммарного количества патентов и программ для ЭВМ.

% Для добавления в список публикаций автора работ, которые не были процитированы в
% автореферате, требуется их~перечислить с использованием команды \verb!\nocite! в
% \verb!Synopsis/content.tex!.
